\chapter{Consumerism and Manufactured Desire}
\label{ch:consumerism}

\epigraph{People recognize themselves in their commodities; they find their soul in their automobile, hi-fi set, split-level home, kitchen equipment.}{Herbert Marcuse, \textit{One-Dimensional Man}}

\noindent
Advertising is the largest doxonomic industry.
Global advertising expenditure surpassed one trillion dollars in 2024, making it the single greatest investment in belief production in human history.
This chapter analyzes consumerism as a doxonomic system.

\section{Advertising as Doxonomic Investment}
\label{sec:ch23-advertising}

\begin{definition}[Advertising as Doxonomic Production]
\label{def:advertising-doxonomic}
\index{advertising}
Advertising is a doxonomic investment whose primary output is not product awareness but \emph{desire infrastructure}: the narrative environment in which consumption appears as the natural expression of identity, freedom, and self-worth.
\end{definition}

The individual advertisement promotes a product.
The advertising \emph{system} promotes a worldview: you are what you buy, happiness is purchasable, identity is consumable \autocite{debord1967, adorno1944}.

\section{The Desire Production Function}
\label{sec:ch23-desire}

\begin{model}[Manufactured Desire]
\label{mod:desire}
\index{manufactured desire}
Let $D_j$ be the desire intensity of agent $j$ for goods beyond subsistence.
Desire is produced by:
\[
  D_j(t) = D_0 + \alpha \sum_k \funding{k}^{\text{ad}} \cdot e_{jk} - \beta \cdot S_j(t)
\]
where $D_0$ is baseline desire, $\funding{k}^{\text{ad}}$ is advertising through channel $k$, $e_{jk}$ is exposure, and $S_j(t)$ is satisfaction with current possessions (which advertising systematically reduces).
The term $-\beta S_j$ captures the mechanism of \emph{induced dissatisfaction} \autocite{schor1998, galbraith1958}.
\end{model}

\section{Consumption as Moral Language}
\label{sec:ch23-moral}

Consumerism is not merely an economic pattern but an anthropological regime: the belief that the good life consists in acquisition, that social worth is legible through possessions, and that consumer choice is the primary form of agency \autocite{fromm1976}.

\section{Notes and References}
\label{sec:ch23-notes}

The critical theory of consumer society draws on \textcite{adorno1944}, \textcite{debord1967}, and \textcite{fromm1976}.
On advertising and identity, see \textcite{veblen1899} and \textcite{galbraith1958}.
Induced dissatisfaction is analyzed in \textcite{schor1998}.
For the economics of manipulation through advertising, see \textcite{akerlof2015}.

\section{Exercises}

\begin{exercise}
Estimate the advertising expenditure per capita in your country and compare it to spending on public education.
Interpret the ratio doxonomically.
\end{exercise}

\begin{exercise}
Using the desire production model, show that if $\beta > 0$, advertising simultaneously increases desire \emph{and} decreases satisfaction, creating a treadmill effect.
\end{exercise}

\begin{exercise}
Analyze a specific advertising campaign as a belief production chain.
Identify funder, institution, channel, narrative, and behavioral consequence.
\end{exercise}
